\documentclass[11pt]{letter}
\usepackage{lmodern}
\usepackage[a4paper,margin=2.1cm]{geometry}
\usepackage[T1]{fontenc}
\usepackage[expansion=false]{microtype}
\setlength{\parskip}{0.4em}
\longindentation=0pt
\indentedwidth=\textwidth
\signature{Prince Solanki, on behalf of the authors\\prince@ma.iitr.ac.in \& prince@iilm.edu}
\address{}
\date{\today}
\begin{document}\pagestyle{empty}
\begin{letter}{The Editors-in-Chief\\Swarm and Evolutionary Computation}
\opening{Dear Editors-in-Chief,}
We present for consideration the manuscript \textit{``Baseline Configuration and Random-Sampling Controls in Metaheuristic Comparisons for Wind Farm Layout Optimization''}, by Prince Solanki, Pulkit Dwivedi, Vanita Garg and Vinay Shukla.

The paper examines how baseline settings, hybrid-component controls and feasibility reporting change comparisons of swarm and evolutionary algorithms. On a 68-case benchmark with 30 paired seeds, it combines evaluation-budget accounting, random-sampling controls, conditional quality rankings and interval-based practical-equivalence assessments. A PSO coefficient sweep connects baseline behaviour to established second-moment stability theory under stagnation; this applies existing theory rather than proposing a new optimizer or convergence theorem.

The revision includes a real-coded GA, re-optimization at $5D$ and $6D$, analytic-gradient SLSQP variants on IEA37 and a second measured-wind block at Lillgrund. GA changes conclusions beyond the original method pool, analytic objective and constraint derivatives strengthen the tested SLSQP implementations, and the SSA control result at $4D$/$5D$ does not persist at $6D$. A new all-run paired reliability analysis complements conditional quality ranks at Lillgrund. These findings make the work relevant to rigorous experimental comparison of swarm and evolutionary methods.

The manuscript distinguishes fixed-benchmark estimates from sensitivities to case exchangeability and grouping. Alternative-model studies transfer final layouts to a different evaluator; the direct spacing and gradient experiments are identified separately. It explains the limits of gradient-cost accounting and the differing constraint-handling implementations.

The accompanying revision archive contains the manuscript, supplementary material, supplied revision records and a standalone archive audit. Original optimization records and several core modules remain to be added for full reproduction, and rounded archival coordinates limit independent strict-feasibility verification. These limitations are disclosed in the manuscript and supplement. Funding, contribution, interest and AI-use statements are provided in the declarations document.

\closing{Sincerely,}
\end{letter}\end{document}
